\chapter{ОНОЛЫН ХҮРЭЭ БА СИСТЕМИЙН ЗАГВАР}

Өмнөх бүлэгт дадал үүсэх онол, зан үйлийг дэмжих дижитал системүүд болон тэдгээрийн хязгаарлалтыг судлагдсан байдлын хүрээнд авч үзсэн. Харин энэ бүлэгт тухайн судалгаануудыг дахин тайлбарлах бус, уг ажлын хүрээнд аль онолын хандлагуудыг сонгон ашиглаж байгаа, мөн тэдгээрийг програм хангамжийн системийн зохиомжид хэрхэн хөрвүүлж буйг тодорхойлно. Өөрөөр хэлбэл, энэ бүлгийн зорилго нь дадал, зан үйлийн онолыг системийн функц, өгөгдлийн бүтэц, алгоритмын логик, харилцан үйлчлэлийн зохиомжтой уялдуулсан нэгдсэн концепц хүрээг боловсруулахад оршино.

Уг судалгаанд онолын тулгуур болгон дадал төлөвшлийн мөчлөг, Фоггийн зан үйлийн загвар, итгүүлэн нөлөөлөх системийн зохиомжийн зарчим, мөн дадлын хүчийг автоматжилтаар хэмжих арга-ыг ашиглана. Эдгээрийг тус тусад нь хэрэглэх бус, харин хэрэглэгчийн дадал төлөвших үйл явцыг үе шаттай дэмжих системийн нэгдсэн логик болгон авч үзсэн.

\section{Судалгаанд ашиглах онолын хандлагын сонголт}

Энэ судалгаанд ашигласан онолын аргууд нь системийн өөр өөр түвшний шийдлийг үндэслэх үүрэгтэй. Дадал төлөвшлийн мөчлөг нь хэрэглэгчийн зан үйлийг өдөөгч дохио, хүсэл тэмүүлэл, хариу үйлдэл, урамшуулал гэсэн үндсэн бүрэлдэхүүнээр задлан авч үзэх боломж олгож байгаа тул системийн үндсэн логик урсгалыг тодорхойлоход ашиглагдана. Фоггийн зан үйлийн загвар нь тухайн зан үйл бодитой хэрэгжих нөхцөлийг сэдэл, боломж, өдөөгч гэсэн гурван хүчин зүйлээр тайлбарладаг учраас мэдэгдлийн логик, төвөгшлийг бууруулах шийдэл, дасан зохицох дэмжлэгийг онолын түвшинд зөвтгөнө. Харин итгүүлэн нөлөөлөх системийн зохиомжийн зарчим нь хэрэглэгчийг албадахгүйгээр, харин дэмжих, хялбаршуулах, ахицыг нь ойлгомжтой болгох зохиомжийн шийдлийг боловсруулахад чиглүүлнэ. Дадлын хүчийг автоматжилтаар хэмжих арга нь системийн үнэлгээ, ахицын шинжилгээ, сануулгыг аажмаар бууруулах шийдвэр болон дасан зохицох логикийн суурь үзүүлэлтийг тогтооно.

\section{Дадал төлөвшлийн мөчлөгийг системийн түвшинд хөрвүүлсэн нь}

Энэ судалгаанд дадал төлөвшлийн мөчлөгийг системийн үндсэн концепц загвар болгон ашиглаж байна. Өөрөөр хэлбэл, системийн гол функцууд нь хэрэглэгчийн зан үйлийг зүгээр нэг бүртгэх бус, харин өдөөгч дохио үүсгэх, сэдлийг дэмжих, гүйцэтгэлийг хялбаршуулах, урамшууллаар бататгах гэсэн дараалсан логикоор зохион байгуулагдана.

\subsection{Өдөөгч дохионы давхарга}

Өдөөгч дохионы давхарга нь хэрэглэгчийн зан үйлийг эхлүүлэхэд чиглэсэн системийн логикийг илэрхийлнэ. Энэ түвшинд систем нь хэрэглэгчийн өдрийн цаг, өмнөх үйл ажиллагаа, сонгосон дадлын хуваарь, мэдэгдлийн хариу үйлдлийн түүх зэрэг мэдээлэлд тулгуурлан тохиромжтой мөчийг тодорхойлох боломжтой. Иймээс өдөөгч дохио нь зөвхөн энгийн сануулга бус, харин хэрэглэгчийн нөхцөл байдалд нийцсэн, цаг хугацааны болон зан үйлийн контекстэд суурилсан дэмжлэг байх шаардлагатай.

\subsection{Хүсэл тэмүүллийн давхарга}

Хүсэл тэмүүллийн давхарга нь хэрэглэгч тухайн зан үйлийг яагаад хийх ёстой вэ гэсэн асуултад хариулах дэмжлэгийн хэсэг юм. Энэ түвшинд систем нь хэрэглэгчийн зорилго, хувийн шалтгаан, өөрийн дүр төрхтэй холбогдсон ойлголт, тухайн дадлаас хүлээгдэж буй ашиг тусыг хадгалж, шаардлагатай үед сануулах байдлаар сэдлийг тогтвортой байлгахад чиглэнэ.

\subsection{Хариу үйлдлийн давхарга}

Хариу үйлдлийн давхарга нь хэрэглэгчийн бодит гүйцэтгэлийг дэмжих логикийг илэрхийлнэ. Энэ түвшинд системийн гол зорилго нь хэрэглэгчийн хийх ёстой алхмыг аль болох энгийн, ойлгомжтой, бага төвөгшилтэй байлгах явдал юм. Иймд нэг товшилтоор бүртгэх, дадлыг жижиг алхамд хуваах, хэсэгчлэн гүйцэтгэлийг зөвшөөрөх, ``хамгийн бага хувилбар'' санал болгох зэрэг шийдлүүд энэ давхаргад хамаарна.

\subsection{Урамшууллын давхарга}

Урамшууллын давхарга нь зан үйлийн дараах бататгал, ахицын мэдрэмж, үргэлжлүүлэх шалтгааныг бэхжүүлэхэд чиглэнэ. Энэ түвшинд систем нь зөвхөн гаднын шагнал өгөхөөс хэтрэхгүй, харин хэрэглэгчийн хийсэн үйлдэл бодит ач холбогдолтой байсан гэдгийг ойлгуулах шаардлагатай. Иймээс системийн урамшууллын логикт богино хугацааны урамшуулал, ахицын өөрчлөлтийн дүрслэл, эргэцүүлсэн санал сэтгэгдэл, ашиг тусын эргэн сануулалт зэрэг механизмууд орно.

\section{Фоггийн зан үйлийн загварыг системийн зохиомжид ашигласан нь}

Фоггийн зан үйлийн загвар нь уг судалгаанд системийн харилцан үйлчлэлийн зохиомж болон дасан зохицох дэмжлэгийн логикийг тодорхойлоход хэрэглэгдэнэ. Тус загварын дагуу зан үйл бодитоор хэрэгжихийн тулд хэрэглэгчид тухайн мөчид хангалттай сэдэл байх, үйлдэл хийх боломж бүрдсэн байх, мөн зөв цагт өдөөгч дохио ирсэн байх шаардлагатай.

Энэ логикийг системийн түвшинд дараах байдлаар тусгаж байна. Нэгдүгээрт, хэрэглэгчийн сэдэл буурах магадлалтай үед зорилгын утга, хувийн шалтгаан, ашиг тусыг сануулах мэдээлэл үзүүлнэ. Хоёрдугаарт, үйлдлийн төвөгшил өндөр байвал гүйцэтгэлийг жижигрүүлэх, нэг товшилтоор хариу өгөх, хэсэгчилсэн гүйцэтгэл зөвшөөрөх зэргээр боломжийг нэмэгдүүлнэ. Гуравдугаарт, мэдэгдлийг тогтмол нэг цагт бус, харин хэрэглэгчийн өдрийн хэв маяг, бүтэлгүйтлийн түүх, санал сэтгэгдэлд тулгуурлан илүү тохиромжтой мөчид хүргэх замаар өдөөгч дохионы үр нөлөөг сайжруулна.

\section{Итгүүлэн нөлөөлөх технологийн зарчмыг системд тусгасан нь}

Итгүүлэн нөлөөлөх технологийн хандлага нь уг судалгаанд хэрэглэгчийг шахах бус, харин дэмжих, ахицыг нь ойлгомжтой болгох, гүйцэтгэлийг хялбаршуулах, эерэг дохио өгөх зохиомжийн зарчмыг тодорхойлоход ашиглагдана. Энэ хүрээнд хэрэглэгчийн ахицыг харуулах, өөрийгөө хянах боломж олгох, богино эерэг бататгал өгөх, төвөгшлийг багасгах, сонголтын эрхийг хадгалах зэрэг нь системийн зохиомжийн үндсэн чиглэл болно.

\section{Дадлын хүчийг тооцоолох нь}

Энэхүү судалгаанд дадлын хүч-ийг зөвхөн гүйцэтгэлийн тоо, тасралтгүй гүйцэтгэлийн цуваа, эсвэл гүйцэтгэлийн хувиар бус, харин тухайн зан үйл хэр зэрэг автоматжсан бэ гэдгээр тодорхойлно. Иймээс уг ажлын хүрээнд дадлын хүчийг хэмжих үндсэн баталгаажсан үзүүлэлтээр SRBAI-д суурилсан автоматжилтын оноо-г сонгов \cite{ref9}.

Хэрэв $q_1, q_2, q_3, q_4$ нь SRBAI-ийн дөрвөн зүйлийн хариулт бол автоматжилтын түүхий оноог дараах байдлаар тооцно.

\[
A_{\text{raw}}=\frac{q_1+q_2+q_3+q_4}{4}
\]

Хэрэв 1--7 Лайкертийн шатлал ашиглавал 0--100 хооронд нормчлогдсон оноог дараах байдлаар тодорхойлно.

\[
A_{100}=\frac{A_{\text{raw}}-1}{6}\times 100
\]

Жишээлбэл, хэрэглэгч тухайн дадалд 5, 6, 5, 4 гэсэн үнэлгээ өгсөн бол:

\[
A_{\text{raw}}=\frac{5+6+5+4}{4}=5.0,\qquad
A_{100}=\frac{5.0-1}{6}\times 100 \approx 66.7
\]

Уг үр дүн нь тухайн дадал дунджаас дээш түвшинд автоматжиж эхэлсэн боловч бүрэн тогтсон гэж үзэхэд эрт үе шатанд байгааг илэрхийлж болно.

Гэсэн хэдий ч зөвхөн өөрийн үнэлгээнд тулгуурлах нь системийн түвшинд хангалтгүй байж болох тул энэхүү судалгаанд гүйцэтгэлийн тогтвортой байдал болон нөхцөл байдлын тогтвортой байдал зэрэг зан үйлийн барагцаалсан үзүүлэлтүүдийг нэмэлтээр тооцно. Үүний үндсэн дээр системийн түвшинд дараах нийлмэл оноог ашиглана.

\[
HS = 0.60\times SRBAI_{100} + 0.25\times Consistency_{100} + 0.15\times ContextStability_{100}
\]

Энд:

\begin{itemize}
    \item $SRBAI_{100}$ нь автоматжилтын үндсэн нормчлогдсон оноо,
    \item $Consistency_{100}$ нь гүйцэтгэсэн боломжит тохиолдлын харьцаа,
    \item $ContextStability_{100}$ нь ижил цагийн хүрээ, өдрийн төрөл, өдөөгч нөхцөлд гүйцэтгэсэн хувийн үзүүлэлт юм.
\end{itemize}

Энэ томьёо нь судалгаанд тогтсон стандарт томьёо биш, харин судалгааны үндэслэлтэй тооцооллын модель болно. Энд SRBAI-д хамгийн өндөр жин өгч буй нь дадлын хүчний хамгийн ойрын шууд хэмжүүр нь автоматжилт байдагтай холбоотой.

Түүнчлэн олон алхамт, урт хугацааны зан үйлийн хувьд бүх үйл явцыг бүхэлд нь автоматждаг гэж үзэхээс илүү тухайн үйлдлийг эхлүүлэх байдал нь автоматждаг гэж тайлбарлах нь онолын хувьд илүү тохиромжтой \cite{ref1}. Иймээс систем нь урт хугацааны дадлын хувьд ``өгөгдсөн өдөөгч дохион дээр эхэлсэн эсэх'' үзүүлэлтийг тусгайлан ажиглана.

Энэхүү судалгаанд дадлын хүчийг дараах ангиллаар тайлбарлаж болно: 0--39 -- сул, 40--69 -- хөгжиж буй, 70--100 -- хүчтэй. Харин ``дадал төлөвшсөн'' гэж системийн түвшинд тодорхойлохдоо дараах нөхцөлийг хамтад нь авч үзнэ: SRBAI-ийн оноо 70 ба түүнээс дээш байх, уг түвшин дор хаяж 2--4 долоо хоног тогтвортой хадгалагдсан байх, нөхцөл байдлын тогтвортой байдал өндөр байх, мөн гүйцэтгэлийн тогтвортой байдал хангалттай түвшинд хүрсэн байх.

\section{Онолын нэгдсэн хүрээ ба системийн концепц загвар}

Энэ судалгаанд ашигласан онолын хандлагуудыг нэгтгэн үзвэл зан үйлд суурилсан дадал төлөвшүүлэх систем нь дараах нэгдсэн логикоор ажиллана. Нэгдүгээрт, систем хэрэглэгчийн нөхцөл байдалд тохирсон өдөөгч дохио үүсгэнэ. Хоёрдугаарт, тухайн дадлын зорилго, хувийн ач холбогдол, хүлээгдэж буй ашиг тусыг сануулснаар хүсэл тэмүүллийг идэвхжүүлнэ. Гуравдугаарт, гүйцэтгэлийг аль болох хялбар болгож, төвөгшлийг бууруулснаар хариу үйлдлийг дэмжинэ. Дөрөвдүгээрт, гүйцэтгэлийн дараа шууд болон эргэцүүлсэн урамшуулал өгч зан үйлийг бататгана. Тавдугаарт, автоматжилтын түвшин болон зан үйлийн прокси үзүүлэлтүүдийгд тулгуурлан дадлын хүчийг үнэлж, дараагийн сануулга, дасан зохицох дэмжлэгийн шийдвэрийг үндэслэнэ.

\section{Бүлгийн дүгнэлт}

Дадал төлөвшлийн мөчлөгийг системийн үндсэн логик, Фоггийн зан үйлийн загварыг харилцан үйлчлэл болон дасан зохицох дэмжлэгийн үндэс, итгүүлэн нөлөөлөх технологийг дэмжих зохиомжийн зарчим, харин дадлын хүчний ойлголтыг автоматжилт ба зан үйлийн прокси үзүүлэлтүүдэд тулгуурласан үнэлгээний хүрээ болгон тусгасан онолын нэгдсэн загвар тодорхойлогдов.
